On considère $n$ jetons alignés, initialement tous blancs.

À chaque opération, on retourne deux jetons voisins.
\begin{enumerate}
  \item Montrer que toute configuration accessible contient un nombre pair de
  jetons noirs.
  \item Soient $i < j$. Montrer qu'en retournant successivement les paires
  \[
  (i, i+1), (i+1, i+2), \ldots, (j-1, j),
  \]
  seuls les jetons situés aux positions $i$ et $j$ changent finalement de
  couleur.
  \item En déduire que toute configuration contenant un nombre pair de jetons
  noirs est accessible.
  \item Conclure qu'une configuration est accessible si et seulement si elle
  contient un nombre pair de jetons noirs.
\end{enumerate}
